\section{Ethical Considerations}
\label{sec:ethics}

Facial affect inference is a dual-use technology, and much of its deployed use
--- inferring the internal states of people who have not asked for it --- is
precisely what our design is meant to resist. We state our position explicitly.

\paragraph{What the system claims.} It estimates a \emph{perceived} affective
display along valence and arousal, as annotated by human raters. It does not
recover a person's felt emotion, intent, or truthfulness, and the dimensional
target is not a diagnosis. The scientific premise that discrete internal
emotional states can be read off a face is contested~\cite{russell1980circumplex};
we therefore predict the continuous, annotator-defined quantity our labels
actually contain, and attach an interval to every prediction rather than a
confident point.

\paragraph{Intended and excluded uses.} The demonstrator estimates the user's own
state, on their own device, at their request. We explicitly exclude monitoring of
people who are not the operator --- employee, student, candidate, or suspect
assessment --- and any deception-detection or mental-health screening
application. Abstention is a first-class outcome for exactly this reason: a
system that must always answer invites use in settings where it should not be
used at all.

\paragraph{Privacy.} No imagery or biometric signal leaves the device in the
demonstrator; inference runs in the browser and no frames are transmitted or
stored. We do not redistribute AffectNet, which is registration-gated, and we
commit no image data to the public repository --- only metrics and rendered
figures.

\paragraph{Bias and unequal performance.} We report aggregate accuracy and
calibration only. AffectNet's demographic composition is not balanced, and
expression annotation is known to carry cultural and racial bias; aggregate
calibration can therefore hide substantially worse calibration for
under-represented groups. We make no fairness claim, and we regard per-group
calibration as the most important extension of this work: an abstention threshold
calibrated on a population that does not represent the user is a threshold that
fails silently for that user, which is the exact failure mode this paper argues
against.
